
同一带电弱流在轻费米子衰变中表现为低能四费米子作用，在重费米子质量足以产生实 $W$ 时则必须回到完整的 $W$ 传播与偏振结构。顶夸克提供了最干净的例子：由于 $m_t>m_W+m_b$，主导衰变 $t\to bW^+$ 是真正的二体过程，而且 $V-A$ 顶角直接决定末态 $W$ 的螺旋度比例。

质量本征态中的带电流顶角给出
\begin{equation}
\boxed{
\mathcal M(t\to bW^+)
=\frac{g}{\sqrt2}V_{tb}\,
\bar u_b(p_b)\gamma^\mu P_Lu_t(p_t)\,
\varepsilon_\mu^*(k).}
\label{eq:sm-top-bw-amplitude-dp94}
\end{equation}
在 $m_b/m_t\to0$ 的受控极限中，令
\begin{equation*}
 x_W=\frac{m_W^2}{m_t^2},
 \qquad
 |\mathbf k|=\frac{m_t c}{2}(1-x_W).
\end{equation*}
对初态顶夸克自旋平均、末态自旋与 $W$ 极化求和后，树级宽度为
\begin{equation}
\boxed{
\Gamma(t\to bW^+)
=\frac{G_F^{(E)}(m_tc^2)^3}{8\pi\sqrt2\,\hbar}
|V_{tb}|^2(1-x_W)^2(1+2x_W)
+\order\!\left(\frac{m_b^2}{m_t^2}\right).}
\label{eq:sm-top-width-dp94}
\end{equation}
这里 $(1-x_W)^2$ 包含二体相空间与振幅中的动量因子，而 $1+2x_W$ 来自纵向和横向 $W$ 极化的总和。

更有信息量的是把宽度按 $W$ 的螺旋度分解。以顶夸克静止系中 $W$ 的飞行方向为量子化轴，在 $m_b=0$ 时
\begin{equation}
\boxed{
F_0=\frac{1}{1+2x_W},
\qquad
F_-=\frac{2x_W}{1+2x_W},
\qquad
F_+=0.}
\label{eq:sm-top-helicity-fractions-dp94}
\end{equation}
$F_+$ 在无质量 $b$ 极限严格为零，不是相空间巧合，而是左手带电流与角动量守恒共同造成的螺旋度选择定则。有限 $m_b$、QCD 辐射与电弱圈修正会使它变成很小但非零的量。纵向分量 $F_0$ 在 $m_t\gg m_W$ 时趋于 $1$，与 Goldstone 等价定理一致：高能纵向 $W$ 的主导自由度来自 Higgs 双重态中的 Goldstone 模式。

\begin{figure}[htbp]
 \centering
 \includegraphics[width=.72\textwidth]{sm_top_w_helicity_dp94.pdf}
 \caption{树级 $t\to bW$ 中 $m_b=0$ 时的 $W$ 螺旋度分数。横轴取 $m_W/m_t$；当质量层级 $m_t\gg m_W$ 增强时，纵向分量 $F_0$ 按 Goldstone 等价定理趋于主导。曲线由\eqnrefs{eq:sm-top-helicity-fractions-dp94} 直接计算。}
 \label{fig:sm-top-w-helicity-dp94}
\end{figure}
